\section{Introduction: multiplication through a finite contact algebra}
\label{sec:conductor-introduction}
Let $V_n=H^0(\Pj^1,\mathcal O(n))$. A proper subseries
$U\subset V_n$ may fail to generate the complete series in higher
symmetric degree even when it has no base points. We study the
\emph{scheme} of this failure: it is defined by the zeroth Fitting
ideal of the cokernel of $\Sym^m U\to V_{mn}$. The scheme remembers
multiplicities which neither its support nor a tangent calculation
determines.

Our first object is a subseries containing all sections vanishing on
a prescribed finite divisor. Let $D$ have length $d$, put
$W_D=H^0(\mathcal O(n)(-D))$, and let
\[
 A\subset E_{n,D}:=H^0(\mathcal O_D(n)),\qquad
 U_A=\operatorname{res}_D^{-1}(A).
\]
We require $A$ to generate $\mathcal O_D(n)$; repeated points of $D$
are allowed. These are entire open relative Grassmannians of
subseries containing $W_D$, not selections of failure points. For
$n\ge2d-1$ we prove an isomorphism of coherent cokernel sheaves
\begin{equation}\label{eq:intro-conductor-cokernel}
 \operatorname{coker}(\Sym^m\mathcal U_A\longrightarrow V_{mn})
 \simeq
 \operatorname{coker}(\Sym^m\mathcal A\longrightarrow E_{mn,D})
 \quad(m\ge2).
\end{equation}
It holds over the moving-divisor parameter space and after arbitrary
base change. Thus all Fitting schemes, including nonreduced fibres,
are transported to a multiplication problem on a finite algebra.
The point is a global surjectivity statement for products containing
$W_D$, not an elimination of an invertible block at a single point.

For pencils, the finite problem has a complete solution. Write
$G_D^\circ\subset\Gr(2,E_{n,D})$ for the generating locus. On a
smooth frame chart choose a basis $a,b$ with $a$ a unit on $D$,
and put $v=b/a$. The cokernel is presented by
$[1,v,\ldots,v^m]$ in the length-$d$ algebra $H^0(\mathcal O_D)$.
For $m\ge d-1$ the failure ideal is independent of $m$.

\begin{theorem}[Contact-pencil failure schemes]\label{thm:contact-main}
Let $d\ge4$, $n\ge2d-1$, and $m\ge d-1$. The maximal-minor
failure scheme on $G_D^\circ$ is a Cartier divisor. If
$D=\sum_i d_i p_i$, its complete primary decomposition is
\begin{equation}\label{eq:intro-contact-primary}
 \mathcal J_m=
 \bigcap_{i:d_i\ge2}\mathcal I_{T_i}^{\binom{d_i}{2}}
 \ \cap\!
 \bigcap_{i<j}\mathcal I_{E_{ij}}^{d_i d_j}.
\end{equation}
Here $T_i$ is the divisor on which $A$ does not separate $2p_i$,
and $E_{ij}$ is the divisor on which it does not separate $p_i+p_j$.
These are distinct smooth irreducible divisors. There are no embedded
associated points. On smooth frame charts the exact local equation
is a weighted product of linear forms; its weights, nilpotency
indices, successive nilradical quotients, tangent cones, and every
multiplication corank are given in
Theorems~\ref{thm:contact-factorization} and~\ref{thm:contact-layers}.
Corollary~\ref{cor:contact-multiplicity-strata} determines all
components and their primary multiplicities over every moving
multiplicity stratum.

Allowing $D$ to move, the failure divisor on the relative generating
Grassmannian is integral. Its normalization is smooth and is the
finite incidence which marks a length-two subdivisor $Z\le D$
not separated by $A$. Theorem~\ref{thm:moving-normalization} gives
this normalization, and Proposition~\ref{prop:moving-singular-test}
gives an exact test for its singular support.
\end{theorem}

These subseries have codimension $d-2$ in $V_n$, which is unbounded
as $d$ varies. In contrast to the hyperplane family, their
scheme structure stabilizes once $m\ge d-1$. It nevertheless
varies sharply with the multiplicities of $D$: for $D=dp$ the
fixed-divisor failure scheme is the $\binom d2$-fold ramification
divisor, whereas for a reduced divisor it is a reduced arrangement
of pair-identification divisors. Both are fibres of the same
integral moving-divisor scheme. The normalization explains why
nonreduced fixed-contact fibres must not be identified with embedded
components of the total family.

The proof has three steps. A base-point-free pencil argument fills
all products vanishing twice on $D$; the first normal quotient then
fills all products vanishing once. This proves
\eqref{eq:intro-conductor-cokernel}. A power-basis determinant in the
finite algebra gives \eqref{eq:intro-contact-primary}, including
all confluent cases. Finally the marked length-two incidence gives
a finite birational resolution of the moving failure divisor.
The power-basis determinant, confluent Vandermonde identity, and
Fitting base change are classical algebraic tools. The assertions
proved here concern their exact relative realization as
multiplication failure schemes, their full primary fibres, and the
normalization across contact collisions.

The later sections retain the quadratic component classification,
its polar residual deformation theory and wall singularities, and
the all-degree hyperplane theorem. The new conductor argument does
not depend on the statistical interpretation or on those later
classification results. Conversely, it does not identify the
contact-subseries locus with the full ambient Grassmannian. The
unrestricted quadratic excess scheme and the embedded primary ideal
of the hyperplane family are different objects. Their previously
proved statements are retained with their exact hypotheses.
Section~\ref{sec:prior-comparison} distinguishes the classical
inputs from the family-specific conclusions and records the
remaining source-access limitation in the comparison with
\cite{Ballico93}. Full contact, realization and statistical
applications are reproduced in the appendices.
